\clearpage
\phantomsection
\chapter*{题目索引}
\addcontentsline{toc}{chapter}{题目索引}
\par\medskip\noindent\textbf{第1章}\par
\BookIndexEntry{cn-ex1-1}{例 1-1　端到端服务从哪层开始 【2009】}
\BookIndexEntry{cn-ex1-2}{例 1-2　相邻节点间的流量控制 【2022】}
\BookIndexEntry{cn-ch1-review}{自测与答案}
\par\medskip\noindent\textbf{第2章}\par
\BookIndexEntry{cn-ex2-1}{例 2-1　由调制状态数求最大速率 【2022】}
\BookIndexEntry{cn-ex2-2}{例 2-2　增加信噪比是否一定提高固定调制的极限 〖2.1-017〗}
\BookIndexEntry{cn-ex2-3}{例 2-3　时延带宽积先转成传播时间 【2023】}
\BookIndexEntry{cn-ex2-4}{例 2-4　虚电路必须预留带宽吗 【2020】}
\BookIndexEntry{cn-ex2-5}{例 2-5　MAC 地址属于物理层接口规范吗 【2018】}
\BookIndexEntry{cn-ch2-review}{自测与答案}
\par\medskip\noindent\textbf{第3章}\par
\BookIndexEntry{cn-ex3-1}{例 3-1　填充按连续1计数，不按字节重置 【2013】}
\BookIndexEntry{cn-ex3-2}{例 3-2　找出满足生成式的接收串 【2023】}
\BookIndexEntry{cn-ex3-3}{例 3-3　先比较所有码字，再谈纠错 【2025】}
\BookIndexEntry{cn-ex3-4}{例 3-4　长确认不能从周期中消失 【2020】}
\BookIndexEntry{cn-ex3-5}{例 3-5　收到哪些确认，不等于只收到哪些数据 【2009】}
\BookIndexEntry{cn-ex3-6}{例 3-6　先求填满链路的窗口，再求序号位数 〖3.4-019〗}
\BookIndexEntry{cn-ex3-7}{例 3-7　CDMA 先按码片分组，再做内积 【2014】}
\BookIndexEntry{cn-ex3-8}{例 3-8　最短帧缩短后要缩短多少电缆 【2009】}
\BookIndexEntry{cn-ex3-9}{例 3-9　收到 CTS 后只计算剩余时间 【2024】}
\BookIndexEntry{cn-ex3-10}{例 3-10　填充前先识别有无标签 〖3.6-008〗}
\BookIndexEntry{cn-ex3-11}{例 3-11　直通交换最早何时能决定出口 【2013】}
\BookIndexEntry{cn-ch3-review}{自测与答案}
\par\medskip\noindent\textbf{第4章}\par
\BookIndexEntry{cn-ex4-1}{例 4-1　逐列做加法，再比较候选 【2021】}
\BookIndexEntry{cn-ex4-2}{例 4-2　为何800 B链路上的分片总长只有796 B 【2021】}
\BookIndexEntry{cn-ex4-3}{例 4-3　非整字节前缀怎样求地址边界 【2023】}
\BookIndexEntry{cn-ex4-4}{例 4-4　``完整划成5个子网''限制了最深分支 【2019】}
\BookIndexEntry{cn-ex4-5}{例 4-5　四个 \texttt{/21} 应缩短几位 【2018】}
\BookIndexEntry{cn-ex4-6}{例 4-6　相差几个第三字节，仍可能在同一子网 〖4.2-043〗}
\BookIndexEntry{cn-ex4-7}{例 4-7　初始 DHCPREQUEST 的源、目的IP 【2025】}
\BookIndexEntry{cn-ex4-8}{例 4-8　历史 ICMP 题要保留协议年代 【2010】}
\BookIndexEntry{cn-ex4-9}{例 4-9　地址位数增加96，不是地址数乘96 【2023】}
\BookIndexEntry{cn-ch4-review}{自测与答案}
\par\medskip\noindent\textbf{第5章}\par
\BookIndexEntry{cn-ex5-1}{例 5-1　谁决定这个 UDP 数据报交给哪个本机服务 【2018】}
\BookIndexEntry{cn-ex5-2}{例 5-2　少做一次回卷，会错在哪一位 【2024】}
\BookIndexEntry{cn-ex5-3}{例 5-3　最后收到的字节与确认号相差1 【2009】}
\BookIndexEntry{cn-ex5-4}{例 5-4　至少关闭时间要从FIN开始画 【2022】}
\BookIndexEntry{cn-ex5-5}{例 5-5　没有拥塞，不代表一定指数增长 【2020】}
\BookIndexEntry{cn-ex5-6}{例 5-6　ACK 只移走已经确认的那一段 【2025】}
\BookIndexEntry{cn-ex5-7}{例 5-7　拥塞窗口变大后，接收窗口仍可成为上限 【2014】}
\BookIndexEntry{cn-ch5-review}{自测与答案}
\par\medskip\noindent\textbf{第6章}\par
\BookIndexEntry{cn-ex6-1}{例 6-1　问的是本地服务器发出多少次查询 【2016】}
\BookIndexEntry{cn-ex6-2}{例 6-2　传递 FTP 命令使用哪条连接 【2009】}
\BookIndexEntry{cn-ex6-3}{例 6-3　一次 POP3 连接可以读几封邮件 【2025】}
\BookIndexEntry{cn-ex6-4}{例 6-4　非持久且不允许并行，要算上 HTML 本身 【2024】}
\BookIndexEntry{cn-ex6-5}{例 6-5　流水线仍不能跨过资源发现依赖 〖6.5-012〗}
\BookIndexEntry{cn-ex6-6}{例 6-6　DNS 与 Web 的时间由哪几项组成 【2020】}
\BookIndexEntry{cn-ex6-7}{例 6-7　收到 HTML 后，拥塞窗口还限制图片分几轮 【2022】}
\BookIndexEntry{cn-ch6-review}{自测与答案}
\clearpage
\phantomsection
\chapter*{方法索引}
\addcontentsline{toc}{chapter}{方法索引}
\par\medskip\noindent\textbf{第1章}\par
\BookIndexEntry{cn-01-01}{网络由什么组成，提供什么能力}
\BookIndexEntry{cn-01-02}{网络分类与拓扑描述的是哪些维度}
\BookIndexEntry{cn-01-03}{发送、传播、排队：怎样确定时延}
\BookIndexEntry{cn-01-04}{分层与封装：对等层怎样交换数据}
\BookIndexEntry{cn-01-05}{协议、服务和接口分别约束谁}
\BookIndexEntry{cn-01-06}{OSI 与 TCP/IP：先看模型，再看端点}
\par\medskip\noindent\textbf{第2章}\par
\BookIndexEntry{cn-02-01}{数据、信号、码元与两种带宽}
\BookIndexEntry{cn-02-02}{奈奎斯特与香农：哪个条件限制了速率}
\BookIndexEntry{cn-02-03}{编码、调制与 PCM 怎样区分}
\BookIndexEntry{cn-02-04}{电路、报文和分组交换如何沿时间线计算}
\BookIndexEntry{cn-02-05}{数据报与虚电路：连接不等于电路}
\BookIndexEntry{cn-02-06}{传输介质：比较受什么条件约束}
\BookIndexEntry{cn-02-07}{物理层接口规定什么，不规定什么}
\BookIndexEntry{cn-02-08}{中继器与集线器：再生比特，不解释地址}
\par\medskip\noindent\textbf{第3章}\par
\BookIndexEntry{cn-03-01}{一跳链路提供什么功能}
\BookIndexEntry{cn-03-02}{怎样从连续比特中划出帧}
\BookIndexEntry{cn-03-03}{奇偶校验与 CRC：检验通过能推出什么}
\BookIndexEntry{cn-03-04}{码距与海明码：怎样从检验位定位单错}
\BookIndexEntry{cn-03-05}{反馈、计时器与窗口分别解决什么}
\BookIndexEntry{cn-03-06}{停止---等待：为什么序号至少有0和1}
\BookIndexEntry{cn-03-07}{GBN：累计确认越过了哪些帧}
\BookIndexEntry{cn-03-08}{SR：缓存失序帧为何仍不能立即滑窗}
\BookIndexEntry{cn-03-09}{频分、时分、波分与码分怎样共享介质}
\BookIndexEntry{cn-03-10}{ALOHA 与 CSMA：监听为何不能消除冲突}
\BookIndexEntry{cn-03-11}{CSMA/CD：最短帧与网络跨度的关系}
\BookIndexEntry{cn-03-12}{CSMA/CA、确认与 NAV：隐藏站等到何时}
\BookIndexEntry{cn-03-13}{令牌传递与轮询：如何给出有界的使用机会}
\BookIndexEntry{cn-03-14}{局域网与 IEEE 802 的层次}
\BookIndexEntry{cn-03-15}{以太网帧、MAC 地址与有效载荷}
\BookIndexEntry{cn-03-16}{802.11：接收者、发送者和原始端点可能不同}
\BookIndexEntry{cn-03-17}{VLAN：广播域如何跨交换机延伸}
\BookIndexEntry{cn-03-18}{广域链路与 PPP：不需要争用不等于没有成帧}
\BookIndexEntry{cn-03-19}{交换机如何学习、过滤和转发}
\par\medskip\noindent\textbf{第4章}\par
\BookIndexEntry{cn-04-01}{路由与转发：先选路径，再逐包查表}
\BookIndexEntry{cn-04-02}{SDN：控制平面与数据平面怎样分工}
\BookIndexEntry{cn-04-03}{拥塞与流量控制：瓶颈究竟在哪里}
\BookIndexEntry{cn-04-04}{静态、动态与层次路由}
\BookIndexEntry{cn-04-05}{距离---向量：邻居报出的距离还要加哪一段}
\BookIndexEntry{cn-04-06}{链路状态：全网拓扑怎样变成最短路径树}
\BookIndexEntry{cn-04-07}{IPv4 首部、MTU 与分片重组}
\BookIndexEntry{cn-04-08}{子网掩码、CIDR 与可分配地址}
\BookIndexEntry{cn-04-09}{路由聚合与最长前缀匹配}
\BookIndexEntry{cn-04-10}{NAT：为什么地址变了，传输层校验也要变}
\BookIndexEntry{cn-04-11}{ARP 与逐跳 MAC：最终目的和当前下一跳分开算}
\BookIndexEntry{cn-04-12}{DHCP：请求一个地址不等于已经拥有它}
\BookIndexEntry{cn-04-13}{ICMP：报告差错不等于替原分组重传}
\BookIndexEntry{cn-04-14}{IPv6：基本首部固定，不代表没有扩展}
\BookIndexEntry{cn-04-15}{自治系统：域内最短与域间策略分开}
\BookIndexEntry{cn-04-16}{RIP：加1更新，16表示不可达}
\BookIndexEntry{cn-04-17}{OSPF：区域怎样限制拓扑细节}
\BookIndexEntry{cn-04-18}{BGP：AS\_PATH 是路径向量，不是链路状态}
\BookIndexEntry{cn-04-19}{IP 组播：分支复制放在哪里}
\BookIndexEntry{cn-04-20}{移动 IP：位置改变后，固定地址怎样仍能被找到}
\BookIndexEntry{cn-04-21}{路由器内部：查表后为何仍可能排队}
\par\medskip\noindent\textbf{第5章}\par
\BookIndexEntry{cn-05-01}{端口与套接字：一个IP为什么能运行多个服务}
\BookIndexEntry{cn-05-02}{UDP：8字节首部保留了哪些边界}
\BookIndexEntry{cn-05-03}{反码校验和：进位为什么要加回低位}
\BookIndexEntry{cn-05-04}{TCP 序号、确认号与控制位怎样一起读}
\BookIndexEntry{cn-05-05}{三次握手与连接释放：状态变化由哪个事件触发}
\BookIndexEntry{cn-05-06}{确认、重传与 RTT 估计}
\BookIndexEntry{cn-05-07}{接收窗口：总额度与剩余额度怎样区分}
\BookIndexEntry{cn-05-08}{慢开始、拥塞避免与 Reno 的丢包反应}
\BookIndexEntry{cn-05-09}{两个窗口与在途数据：本次究竟还能发多少}
\par\medskip\noindent\textbf{第6章}\par
\BookIndexEntry{cn-06-01}{C/S 与 P2P：角色由谁发起服务请求决定}
\BookIndexEntry{cn-06-02}{DNS 层次、服务器与区}
\BookIndexEntry{cn-06-03}{递归、迭代与缓存：谁又发了一次查询}
\BookIndexEntry{cn-06-04}{FTP：控制命令与文件数据为什么分两条连接}
\BookIndexEntry{cn-06-05}{邮件系统、邮件格式与 MIME}
\BookIndexEntry{cn-06-06}{SMTP 负责发送，POP3 负责读取}
\BookIndexEntry{cn-06-07}{WWW、URL 与资源位置}
\BookIndexEntry{cn-06-08}{HTTP 请求、连接复用、缓存与时延}
